\documentclass[12pt]{article}

\usepackage{sbc-template}
\usepackage{graphicx,url}
\usepackage{float} 
\usepackage[brazil]{babel}
\usepackage{tabularx}
\newcolumntype{C}{>{\centering\arraybackslash}X} % <-- modified
\usepackage[utf8]{inputenc}  
\usepackage[inline]{enumitem}
\usepackage{tcolorbox}
\usepackage{placeins}

%%%%%%%%%%%%%%%%% Pacotes para inserção de algoritmos - BEGIN
\usepackage[ruled,lined]{algorithm2e}

% Pacote para a definição de novas cores
\usepackage{xcolor}
% Definindo novas cores
\definecolor{verde}{rgb}{0.25,0.5,0.35}
\definecolor{jpurple}{rgb}{0.5,0,0.35}
\definecolor{darkgreen}{rgb}{0.0, 0.2, 0.13}
%\definecolor{oldmauve}{rgb}{0.4, 0.19, 0.28}

\usepackage{listings}
\lstset{
    language=Java,
    basicstyle=\ttfamily\small,
    keywordstyle=\color{jpurple}\bfseries,
    stringstyle=\color{red},
    commentstyle=\color{verde},
    morecomment=[s][\color{blue}]{/**}{*/},
    extendedchars=true,
    showspaces=false,
    showstringspaces=false,
    numbers=left,
    numberstyle=\tiny,
    breaklines=true,
    backgroundcolor=\color{white!10},
    breakautoindent=true,
    captionpos=b,
    xleftmargin=0pt,
    tabsize=2
}
%%%%%%%%%%%%%%%%% Pacotes para inserção de algoritmos - END

\sloppy

\title{Avaliação da Viabilidade da Inteligência Artificial na Geração de Templates de Infraestrutura como Código}

\author{David L. Ho\inst{1}}


\address{Instituto de Ciências Exatas e Informática -- Pontifícia Universidade Católica de Minas Gerais
  (PUC)\\
  30.140-100 -- Minas Gerais -- MG -- Brasil
}

\begin{document} 
\raggedbottom

\maketitle
\begin{abstract}
Infrastructure as Code (IaC) has become a widely adopted practice for automating the provisioning and management of cloud resources. At the same time, advances in Artificial Intelligence have enabled the automatic generation of infrastructure artifacts, potentially reducing the effort required to create and maintain them. However, questions remain regarding the quality, compliance, and security of AI-generated artifacts when compared to manually developed solutions. This work presents an experimental evaluation of 240 IaC templates developed manually, generated by Artificial Intelligence, and generated by Artificial Intelligence with organizational context, using automated pipelines based on Terraform, LocalStack, Checkov, and Trivy. The AI scenarios produced a complete three-resource template set in 2.7 to 3.3 minutes, against 77.5 minutes of net manual development time. Functional execution with LocalStack succeeded in 95.6\% of templates generated with organizational context, against 91.1\% without context and 78.3\% manually. In Security Groups, organizational context reduced unrestricted egress from 93.3\% to 26.7\%. The results provide evidence regarding the benefits and limitations of Artificial Intelligence in Infrastructure as Code generation, showing that this adoption is worthwhile in DevOps environments provided it is accompanied by explicit organizational context.
\end{abstract}
\begin{resumo}
A Infraestrutura como Código (IaC) tornou-se uma prática amplamente adotada para automatizar o provisionamento e o gerenciamento de recursos em ambientes computacionais. Paralelamente, o avanço de modelos de Inteligência Artificial tem possibilitado a geração automática de artefatos de infraestrutura, reduzindo o esforço necessário para sua construção. Entretanto, ainda existem dúvidas sobre a qualidade, a conformidade e a segurança desses artefatos quando comparados ao desenvolvimento manual. Este trabalho apresenta uma avaliação experimental de 240 \textit{templates} de IaC desenvolvidos manualmente, gerados por Inteligência Artificial e gerados por Inteligência Artificial com contexto organizacional, usando \textit{pipelines} automatizados baseados em Terraform, LocalStack, Checkov e Trivy. Os cenários de IA produziram o conjunto completo de três recursos em 2,7 a 3,3 minutos, contra 77,5 minutos de desenvolvimento manual líquido. A execução funcional com LocalStack teve sucesso em 95,6\% dos \textit{templates} gerados com contexto organizacional, contra 91,1\% sem contexto e 78,3\% no manual. No Security Group, o contexto organizacional reduziu a saída irrestrita de 93,3\% para 26,7\%. Os resultados permitem compreender os benefícios e limitações do uso de Inteligência Artificial na geração de artefatos de infraestrutura, mostrando que essa adoção vale a pena em ambientes DevOps desde que acompanhada de contexto organizacional explícito.
\end{resumo}

\begin{tcolorbox}
\footnotesize
\textbf{Bacharelado em Engenharia de \emph{Software} - PUC Minas\\
Trabalho de Conclusão de Curso (TCC)} \\

\indent Orientador de conteúdo (TCC I): Danilo Maia - danilomaia@pucminas.br \\
Orientador de conteúdo (TCC I): João Pedro Batisteli - joaobatisteli@pucminas.br\\
Orientador de conteúdo (TCC I): Leonardo Vilela - leonardocardoso@pucminas.br \\
Orientador acadêmico (TCC I): Cleiton Tavares - cleitontavares@pucminas.br \\
Orientador do TCC II: João Caram Santos de Oliveira - caram@pucminas.br\\ \\
Belo Horizonte, 1 de novembro de 2026.
\end{tcolorbox}


\section{Introdução}
\label{sec:introducao}

A computação em nuvem consolidou-se como uma das principais estratégias adotadas pelas organizações para modernizar infraestrutura, ganhar escala e otimizar recursos computacionais. Nesse cenário, a cultura \textit{DevOps} ganhou destaque ao integrar desenvolvimento e operações e incentivar automação, integração contínua e entrega contínua \cite{kim2016devops}, e a \textit{Infrastructure as Code} (\textit{IaC}) se encaixa nesse contexto ao permitir que recursos de infraestrutura sejam definidos e gerenciados por código versionado, reutilizável e passível de automação \cite{morris2016infrastructure}. Paralelamente, o avanço de ferramentas de Inteligência Artificial generativa, capazes de transformar descrições em linguagem natural em código \cite{chen2021codex}, passou a ser explorado também na criação de artefatos de infraestrutura, como alternativa para reduzir o esforço técnico exigido pela \textit{IaC}.

Apesar desse potencial, reduzir o esforço de desenvolvimento não garante que os \textit{templates} produzidos por IA atendam aos requisitos de padronização, conformidade e segurança esperados no desenvolvimento manual, sobretudo quando as regras internas da organização não são fornecidas ao modelo durante a geração. \textbf{O problema investigado neste trabalho é, portanto, compreender como o uso de Inteligência Artificial impacta a eficiência operacional, a padronização e a segurança de \textit{templates} de \textit{IaC}, quando comparado ao desenvolvimento manual e à geração com contexto organizacional.}

Essa questão é relevante porque incorporar novos recursos de nuvem não pode comprometer requisitos técnicos e de segurança. Artefatos de infraestrutura podem apresentar permissões excessivas, exposição de dados e outros problemas de configuração \cite{rahman2021securitysmells,verdet2025security}, além de defeitos de reconciliação do estado da infraestrutura \cite{hassan2024state}, e processos dependentes de atividades manuais podem aumentar o tempo das entregas \cite{forsgren2018accelerate}, o que torna relevante avaliar alternativas que reduzam o esforço de criação de \textit{templates} sem comprometer sua qualidade.

\textbf{O objetivo geral deste trabalho é avaliar o uso da Inteligência Artificial na geração de \textit{templates} de Infraestrutura como Código, considerando eficiência operacional, padronização e segurança.} Para alcançá-lo, foram definidos três objetivos específicos: (i) analisar se a IA reduz o tempo e o esforço necessários para a criação dos artefatos, (ii) avaliar a conformidade dos \textit{templates} gerados com boas práticas e padrões de \textit{IaC}, e (iii) investigar vulnerabilidades de segurança produzidas pelas diferentes abordagens de geração.

Com essa avaliação, buscou-se produzir evidências empíricas sobre os benefícios e as limitações do uso de Inteligência Artificial na geração de artefatos de infraestrutura, em particular sobre até que ponto o fornecimento de contexto organizacional aproxima os \textit{templates} gerados por IA dos padrões de qualidade obtidos no desenvolvimento manual. Essas evidências contribuem para a adoção mais segura e eficiente dessas tecnologias em ambientes \textit{DevOps}.

O restante deste artigo está organizado da seguinte forma. A Seção~\ref{sec:fundamentacao} apresenta os conceitos de \textit{DevOps}, \textit{IaC}, Inteligência Artificial generativa e \textit{CI/CD} que fundamentam o trabalho. A Seção~\ref{sec:trabalhos_relacionados} discute estudos relacionados à qualidade, segurança, testes e geração automatizada de artefatos \textit{IaC}. A Seção~\ref{sec:metodologia} descreve a metodologia adotada para geração, validação e comparação dos \textit{templates}, a Seção~\ref{sec:resultados} apresenta os resultados obtidos, a Seção~\ref{sec:ameacas} discute as ameaças à validade e a Seção~\ref{sec:conclusao} apresenta as considerações finais e os trabalhos futuros.

\section{Fundamentação Teórica}
\label{sec:fundamentacao}

Esta seção apresenta os conceitos de \textit{DevOps}, Infraestrutura como Código, Inteligência Artificial generativa e pipelines de \textit{CI/CD} que sustentam a geração, a validação e a comparação dos artefatos avaliados neste trabalho.

\subsection{\textit{DevOps} e Infraestrutura como Código}

\textit{DevOps} integra desenvolvimento e operações para aumentar a velocidade e a confiabilidade das entregas. Para Kim et al. \cite{kim2016devops}, essa cultura se sustenta nos Três Caminhos (\textit{The Three Ways}): fluxo de trabalho, ciclos de \textit{feedback} e aprendizado contínuo. Em ambientes \textit{cloud}, a automação reduz atividades manuais e permite incorporar alterações de forma mais frequente e controlada.

A Infraestrutura como Código (\textit{IaC}) leva esse princípio à infraestrutura: recursos computacionais passam a ser definidos em arquivos versionáveis, processados por ferramentas como Terraform, AWS CloudFormation e OpenTofu, e submetidos a práticas de engenharia de software como versionamento, revisão de código e automação de mudanças \cite{morris2016infrastructure}. Como qualquer artefato de software, porém, um \textit{template} de \textit{IaC} pode conter erros, permissões excessivas ou divergências entre o estado esperado e o real da infraestrutura, de modo que sua validade sintática não garante, por si só, que o recurso provisionado esteja correto ou seguro. Por isso, revisão e validação automatizada são mecanismos importantes antes da aplicação em ambientes produtivos.

\subsection{Inteligência Artificial Generativa e Modelos de Linguagem}

\textit{Large Language Models} (\textit{LLMs}) são modelos baseados na arquitetura \textit{Transformer} \cite{vaswani2017attention}, capazes de transformar descrições em linguagem natural em código, ainda que sem garantia de precisão \cite{chen2021codex}. Neste trabalho, essa capacidade é usada para gerar arquivos Terraform a partir de requisitos funcionais, e a saída depende do \textit{prompt} e das informações fornecidas, como padrões de nomenclatura, requisitos de segurança e restrições de implementação, que variam entre organizações.

Dois riscos motivam a validação adotada neste experimento. Primeiro, os LLMs podem alucinar, ou seja, produzir informações incorretas de forma aparentemente plausível \cite{ji2023hallucination}. Na geração de código, isso pode resultar em propriedades inexistentes ou atendimento apenas parcial aos requisitos, e diferentes execuções da mesma solicitação podem produzir respostas distintas. Segundo, uma resposta sintaticamente correta não é necessariamente segura: Pearce et al. \cite{pearce2022copilot} identificaram vulnerabilidades em código produzido com apoio de modelos generativos, e em IaC isso pode significar permissões excessivas ou exposição pública sem que o \textit{template} deixe de ser válido. Por isso, neste trabalho, os artefatos gerados sem e com contexto organizacional são submetidos às mesmas verificações aplicadas aos \textit{templates} manuais.

\subsection{Pipelines de Integração e Entrega Contínua}

Pipelines de \textit{CI/CD} automatizam validação, testes e disponibilização de software e infraestrutura, aumentando a repetibilidade das operações e o \textit{feedback} sobre as alterações realizadas \cite{humble2010continuous}. Em \textit{IaC}, isso permite executar de forma padronizada verificações de sintaxe, conformidade e segurança, reduzindo a dependência de checagens manuais. É esse mecanismo que, neste trabalho, permite submeter \textit{templates} manuais e gerados por IA aos mesmos controles, produzindo métricas comparáveis entre os três cenários experimentais.

\section{Trabalhos Relacionados}
\label{sec:trabalhos_relacionados}

Esta seção apresenta estudos empíricos relacionados à segurança, testes, qualidade e geração automatizada de artefatos de Infraestrutura como Código, destacando o método, os principais resultados e a relação de cada um com a proposta desta pesquisa.

Rahman et al. \cite{rahman2021securitysmells} investigaram \textit{security smells} em scripts Ansible e Chef com a ferramenta SLAC, identificando 46.600 ocorrências, incluindo 7.849 senhas codificadas diretamente nos scripts, evidência de que o uso de \textit{IaC} não elimina problemas de segurança. Verdet et al. \cite{verdet2025security} analisaram 812 projetos Terraform \textit{open source} com Checkov e Tfsec, observando maior adoção de políticas de controle de acesso do que de criptografia em repouso, o que reforça que a existência de recomendações de segurança não garante sua aplicação. Diferentemente desses trabalhos, esta pesquisa não analisa apenas artefatos existentes, mas compara \textit{templates} produzidos manualmente e por IA em duas condições de contexto, usando análise estática em arquivos Terraform como um dos instrumentos de avaliação.

Hassan et al. \cite{hassan2024state} estudaram defeitos de reconciliação de estado em \textit{IaC} (divergência entre o estado real e o esperado), identificando oito categorias de falhas com apoio de modelos de linguagem. Os resultados mostram que a ausência de erros sintáticos não garante que o estado produzido corresponda ao esperado, o que sustenta a inclusão, neste trabalho, de verificações de execução além das análises estáticas. Sokolowski, Spielmann e Salvaneschi \cite{sokolowski2024automated} propuseram a \textit{Automated Configuration Testing} para programas de \textit{IaC} de propósito geral, avaliando 6.081 programas entre projetos do GitHub e \textit{benchmarks} artificiais. Enquanto os autores propõem uma técnica específica de teste, esta pesquisa utiliza pipelines de \textit{CI/CD} para aplicar um conjunto comum de verificações a artefatos produzidos por diferentes métodos.

Zhang et al. \cite{zhang2024generative} avaliaram 98 manifestos Kubernetes gerados pelo ChatGPT e identificaram pelo menos um \textit{smell} em 35,8\% deles, sobretudo ausência de limites de CPU e memória, concluindo que modelos generativos produzem configurações funcionais que ainda exigem garantia de qualidade. Esta pesquisa amplia essa perspectiva ao utilizar Terraform e comparar geração manual, IA sem contexto organizacional e IA com contexto organizacional.

Em conjunto, os estudos mostram que artefatos de infraestrutura estão sujeitos a vulnerabilidades, defeitos de estado e problemas de qualidade, e que mecanismos automatizados são relevantes para sua avaliação, cada qual com um foco distinto, seja segurança, defeitos de estado, técnicas de teste ou geração por IA. Nenhum deles, porém, combina em um mesmo experimento desenvolvimento manual, geração por IA sem contexto e geração por IA com contexto organizacional, avaliando simultaneamente eficiência operacional, conformidade, padronização e segurança, lacuna que esta pesquisa busca preencher.

\section{Materiais e Métodos}
\label{sec:metodologia}

Esta pesquisa é de natureza aplicada, com caráter exploratório e descritivo, e adota abordagem predominantemente quantitativa, complementada pela análise qualitativa dos problemas identificados, combinação adequada para comparar, de forma objetiva, tratamentos cujos efeitos ainda não são bem compreendidos na literatura. Do ponto de vista dos procedimentos técnicos, consiste em uma avaliação experimental comparativa, na qual diferentes métodos de geração de \textit{templates} são submetidos à mesma infraestrutura de validação, o que permite comparar os resultados entre os métodos, embora fatores externos ao método de geração, discutidos na Seção~\ref{sec:ameacas}, também possam influenciar parte dos resultados.

A Figura~\ref{fig:fluxo_experimental} apresenta as principais etapas do experimento: definição dos recursos e requisitos, produção dos \textit{templates} nos cenários manual, IA sem contexto organizacional e IA com contexto organizacional, submissão ao mesmo \textit{pipeline} de validação automatizada (validação Terraform, execução funcional com LocalStack, validação de segurança e métricas operacionais) e consolidação dos resultados para comparação.

\begin{figure}[H]
\centering
\includegraphics[width=0.88\textwidth]{diagrama-metodologia.png}
\caption{Fluxo experimental para geração, validação e avaliação de \textit{templates} de Infraestrutura como Código. Fonte: elaborado pelo autor.}
\label{fig:fluxo_experimental}
\end{figure}

\subsection{Ambiente Experimental}

Os artefatos foram desenvolvidos em Terraform, ferramenta de \textit{IaC} \textit{open source} amplamente adotada e compatível com múltiplos provedores de nuvem, e versionados no GitHub, onde \textit{workflows} em GitHub Actions executam formatação, validação, conformidade e segurança (Checkov e Trivy, ferramentas \textit{open source} consolidadas de análise estática) e execução funcional com LocalStack, escolhido por permitir criar, consultar e remover os recursos sem o custo e o risco de uma conta real da AWS. Programas auxiliares em Go e \textit{Shell} apoiam a geração dos artefatos de IA, a organização dos dados e a condução dos testes funcionais.

A geração automatizada foi realizada pelo Claude Code, executado em modo não interativo a partir de um programa em Go, com um único modelo e os mesmos parâmetros em todas as execuções. Os \textit{prompts} base permaneceram inalterados durante a coleta, e a única diferença entre os cenários automatizados é a inclusão ou ausência do contexto organizacional. O modelo, a data das execuções e os arquivos de entrada foram registrados, e as respostas foram armazenadas individualmente no repositório, o que aumenta a rastreabilidade do experimento, ainda que o modelo esteja identificado apenas pelo nome da variante (Claude Sonnet), sem o número exato da versão.

Os \textit{prompts} seguem os elementos Papel, Contexto, Ação, Formato de Saída e Restrições e Avisos. O \textit{prompt} sem contexto omite apenas o contexto organizacional em relação ao prompt com contexto, para não induzir uma divergência artificial entre os dois cenários automatizados.

\subsection{Estratégia de Amostragem}

Foram avaliados \textit{templates} referentes a Amazon S3, AWS IAM e Security Group, escolhidos por representarem armazenamento, gerenciamento de identidades e controle de acesso em rede.

Foram recrutados 20 participantes entre estudantes e profissionais de tecnologia, sem exigência de experiência prévia em Terraform ou \textit{IaC}. O tamanho da amostra manual foi limitado pela dificuldade de recrutar voluntários que concluíssem integralmente a atividade durante o período de coleta, discutida na Seção~\ref{sec:ameacas}. Um formulário de caracterização registrou vínculo acadêmico ou profissional, experiência prévia com nuvem, \textit{IaC} e Terraform, e familiaridade com os serviços avaliados, para descrever o perfil da amostra e auxiliar na interpretação de diferenças de desempenho. Um formulário de registro de atividade, respondido ao final da tarefa, registrou horário de início e término, pausas, dificuldade percebida e eventuais dúvidas sobre os requisitos. Foram incluídos participantes que concluíram integralmente as atividades e forneceram os registros necessários de tempo e submissão.

Cada participante produziu manualmente um \textit{template} por recurso, totalizando 60 artefatos manuais. Para os cenários de IA, cada recurso foi gerado 30 vezes sem contexto e 30 vezes com contexto, totalizando 180 \textit{templates} automatizados. Ao todo, este trabalho avalia 240 \textit{templates}: 180 automatizados e 60 manuais, todos com a mesma especificação funcional, para preservar condições equivalentes de comparação.

\subsection{Contexto Organizacional}

O contexto organizacional foi definido a partir de diretrizes elaboradas para o experimento: padrões de nomenclatura, \textit{tags} obrigatórias com chaves específicas, nomenclatura fixa de variáveis e \textit{outputs}, estrutura obrigatória do README, e versões mais específicas dos requisitos mínimos de segurança já exigidos na linha de base dos três cenários: por exemplo, as quatro \textit{flags} do \textit{Public Access Block} no S3 em vez de apenas ``bloquear acesso público'', ou a proibição de \texttt{0.0.0.0/0} fora da porta 443/tcp no Security Group em vez de apenas ``evitar abertura irrestrita''.

Essas regras representam práticas comuns em organizações que adotam \textit{IaC}, mas que não fazem parte da descrição funcional de um recurso. O objetivo é verificar se tornar essas práticas mais específicas, somadas às convenções de nomenclatura e documentação, influencia a conformidade, a padronização e a segurança dos artefatos gerados. Todo o material está armazenado no repositório e permanece inalterado durante a coleta, variando apenas sua presença entre os dois cenários automatizados.

\subsection{Desenvolvimento do Ambiente Experimental}

\subsubsection{Organização do Repositório}

O ambiente foi implementado no repositório público disponível em \url{https://github.com/Davidlh80/tcc}, permitindo que participantes externos acessem os requisitos e submetam os \textit{templates} manuais por \textit{Pull Request}, na mesma esteira de validação usada nos cenários de IA. O repositório centraliza requisitos, \textit{prompts}, contexto organizacional, código Terraform, scripts, pipelines e resultados, com os documentos de entrada mantidos separados dos artefatos gerados, de modo a evitar que informações exclusivas do cenário com contexto fossem usadas no cenário sem contexto.

Os \textit{templates} foram organizados nos diretórios \texttt{manual}, \texttt{ia-sem-contexto} e \texttt{ia-com-contexto}. No cenário manual, cada participante recebeu um identificador entre \texttt{p01} e \texttt{p30}, usado apenas para organizar as submissões durante o \textit{merge} para a \textit{branch} principal. Nos automatizados, cada geração recebe um identificador entre \texttt{exec-01} e \texttt{exec-30}, cada um com diretórios próprios para S3, IAM e Security Group, contendo \texttt{main.tf}, \texttt{variables.tf}, \texttt{outputs.tf}, \texttt{versions.tf} e \texttt{README.md}. Nos cenários de IA, o programa também cria um \texttt{metadata.yml} com cenário, execução, recurso, modelo e arquivos de entrada, para verificar automaticamente a presença dos componentes esperados e rastrear a origem dos artefatos.

\subsubsection{Preparação da Atividade Manual e Geração por IA}

Cada participante recebeu um identificador, acesso ao repositório, uma \textit{branch} \texttt{manual/pXX} e os requisitos gerais e específicos, devendo alterar somente os três diretórios associados ao seu identificador. Os requisitos descrevem o comportamento esperado: bloqueio de acesso público, criptografia, versionamento e \textit{tags} no S3, uma \textit{policy} de menor privilégio anexada a uma \textit{role} com \textit{trust policy} restrita no IAM, e regras explícitas de entrada e saída no Security Group. Eles não forneceram, porém, implementação pronta, e o participante não recebeu credenciais AWS. Início, término, pausas, dificuldade percebida e experiência prévia foram registrados em formulário, e a primeira submissão foi considerada definitiva.

Para os dois cenários de IA, um programa em Go lê os arquivos de entrada, invoca o Claude Code em modo não interativo e armazena a resposta no diretório correspondente. As até três tentativas automáticas previstas pelo programa cobrem apenas falhas técnicas da chamada em si, como \textit{timeout} ou erro de processo, e não avaliam o conteúdo da resposta. No cenário sem contexto, a entrada contém somente o \textit{prompt} do recurso. Com contexto, o programa acrescenta o documento de padrões organizacionais ao mesmo requisito técnico, mantendo o mesmo modelo nos dois casos. Foram realizadas 30 execuções dos três recursos em cada cenário (90 \textit{templates} sem contexto, 90 com contexto), e se algum dos cinco arquivos previstos não for identificado na resposta recebida, a execução é registrada como falha, sem nova tentativa e sem correção manual dos artefatos em nenhum dos três cenários antes da avaliação.

\subsubsection{Submissão e Fluxo dos Pipelines}

Nos três cenários, os artefatos são submetidos ao mesmo repositório e passam pelos mesmos quatro \textit{workflows} do GitHub Actions, acionados por alteração nos arquivos Terraform, por \textit{Pull Request} ou manualmente, com exceção do \texttt{localstack.yml}, executado apenas sob demanda. Para preservar a equivalência, é considerada para avaliação a primeira versão completa submetida, sem correções posteriores: o retorno dos pipelines funciona como instrumento de avaliação, não como mecanismo de aperfeiçoamento de apenas um dos métodos. A Tabela~\ref{tab:workflows} resume as etapas e a saída de cada \textit{workflow}.

\begin{table}[H]
\centering
\caption{Workflows de validação automatizada. Fonte: elaborado pelo autor.}
\label{tab:workflows}
\small
\begin{tabularx}{\textwidth}{>{\raggedright\arraybackslash}p{2.8cm}|>{\raggedright\arraybackslash}X|>{\raggedright\arraybackslash}p{2.1cm}}
\hline
\textbf{Workflow} & \textbf{Etapas} & \textbf{Saída} \\
\hline
\texttt{terraform-}\\\texttt{validate.yml} & \texttt{fmt -check}, \texttt{init} (sem \textit{backend}), \texttt{validate} & CSV, JSON e MD \\
\hline
\texttt{security-}\\\texttt{scan.yml} & Checkov e Trivy, em paralelo e sem bloquear o \textit{workflow} & JSON \\
\hline
\texttt{metrics.yml} & inventário dos arquivos esperados por \textit{template} & JSON e CSV \\
\hline
\texttt{localstack}\\\texttt{.yml} & \texttt{init}, \texttt{plan}, \texttt{apply}, verificação do recurso, \texttt{destroy} (até 3 lotes em paralelo) & CSV e \textit{logs} \\
\hline
\end{tabularx}
\end{table}

No \texttt{security-scan.yml}, os achados do Checkov e do Trivy são os próprios dados da pesquisa e por isso não interrompem o \textit{workflow}: uma execução concluída indica apenas que a análise foi realizada, não que o \textit{template} está livre de problemas. No \texttt{localstack.yml}, cada lote sobe um contêiner com suporte a S3, IAM, EC2 e STS, e o script \texttt{run-localstack-tests.sh} executa \texttt{init}, \texttt{plan} e \texttt{apply}, consulta via AWS CLI o recurso criado e então o remove (\texttt{destroy}). Cada etapa é registrada separadamente em CSV, e uma falha não interrompe o restante das execuções.

\subsection{Variáveis, Coleta e Métricas}

A variável independente corresponde ao método de geração do \textit{template}: manual, IA sem contexto ou IA com contexto organizacional. As variáveis dependentes, isto é, tempo de desenvolvimento, percentual de conformidade, quantidade e severidade de problemas de segurança, falhas nas validações automatizadas e taxa de sucesso de execução, correspondem diretamente aos três objetivos específicos da pesquisa (esforço, conformidade e segurança), respondendo a cada um com uma métrica própria.

Os principais instrumentos de coleta são o formulário de caracterização dos participantes, os relatórios dos \textit{workflows} de CI/CD e os arquivos \texttt{metadata.yml} do programa de geração. Nas atividades manuais, cada participante registrou em formulário horário de início e término, pausas, dificuldade percebida e experiência prévia, com os registros do GitHub como apoio de auditoria. Nas gerações por IA, o programa registrou os instantes de início e término e o tempo total para disponibilizar os arquivos, considerado na comparação principal porque representa o intervalo até o artefato estar pronto para submissão.

A qualidade técnica foi avaliada pelos resultados de \texttt{terraform fmt}, \texttt{init}, \texttt{validate} e pela execução funcional com LocalStack, observando separadamente \texttt{plan}, \texttt{apply}, verificação do recurso e \texttt{destroy}. Checkov e Trivy forneceram os achados de segurança, incluindo a severidade das configurações inseguras. O percentual de conformidade corresponde à razão entre as verificações do Checkov aprovadas e o total de verificações aplicáveis, agregadas por recurso e cenário, e a padronização foi avaliada pelo inventário do \texttt{metrics.yml}, que verifica a presença dos arquivos esperados em cada \textit{template}. A taxa de sucesso funcional corresponde à proporção de artefatos que concluíram todas as etapas do LocalStack, permitindo identificar em qual etapa cada falha ocorreu.

\subsection{Análise dos Resultados}

Os dados foram analisados por estatística descritiva (médias, medianas e frequências absolutas e relativas), uma vez que o interesse está em comparar diretamente as proporções obtidas por cada tratamento, sem assumir que os dados sigam alguma distribuição estatística específica (como a distribuição normal).

Essa análise quantitativa foi complementada pelo exame qualitativo dos principais problemas encontrados nos artefatos, especialmente quando uma mesma categoria de falha se repetiu em determinado cenário. Essa combinação permite discutir não apenas a quantidade de problemas, mas também sua natureza e seu possível impacto sobre a adoção de IA na geração de Infraestrutura como Código.

\section{Resultados}
\label{sec:resultados}

Esta seção apresenta os resultados dos três cenários avaliados. Todos os dados provêm de uma única execução de referência dos \textit{workflows} descritos na Seção~\ref{sec:metodologia}, realizada sobre a \textit{branch} principal do repositório, com Terraform 1.9.8, \textit{provider} AWS 5.100.0 na configuração padrão (exceto três \textit{templates} manuais executados com a versão 6.68.0, conforme discutido na Seção~\ref{sec:ameacas}), Checkov 3.4.0 e Trivy 0.70.0. Foram avaliados os 180 \textit{templates} automatizados (30 execuções dos três recursos em cada cenário de IA) e os 60 \textit{templates} submetidos pelos 20 participantes do cenário manual (identificadores \texttt{p01} a \texttt{p30}). O experimento compreende, portanto, 240 \textit{templates}.

\subsection{Eficiência Operacional}

A Tabela~\ref{tab:tempo} apresenta o tempo total de geração automatizada (do início da chamada ao Claude Code até a resposta completa), em segundos, registrado no campo \texttt{total\_duration\_ms} dos arquivos \texttt{metadata.yml} das 180 execuções automatizadas, e o tempo líquido de desenvolvimento manual, em minutos, registrado pelos 20 participantes no formulário de registro de atividade (horário de início e término, descontadas as pausas declaradas). Para cada identificador de execução, os tempos dos três recursos (S3, IAM e Security Group) foram agregados em uma única observação, produzindo 30 observações por cenário de IA e 20 observações no cenário manual.

\begin{table}[H]
\centering
\caption{Tempo de produção do conjunto dos três \textit{templates} (S3, IAM e Security Group), por cenário. Fonte: elaborado pelo autor.}
\label{tab:tempo}
\small
\begin{tabular}{l|c|c}
\hline
\textbf{Cenário} & \textbf{Média} & \textbf{Mediana} \\
\hline
Manual (n=20, minutos) & 77,5 min & 80,0 min \\
IA sem contexto (n=30, segundos) & 160,8 s (2,7 min) & 150,0 s \\
IA com contexto (n=30, segundos) & 195,9 s (3,3 min) & 208,5 s \\
\hline
\end{tabular}
\end{table}

A geração automatizada do conjunto dos três \textit{templates} leva em média 2,7 minutos sem contexto organizacional e 3,3 minutos com contexto, contra 77,5 minutos de desenvolvimento manual líquido, excluídas as pausas declaradas. O formulário de registro de atividade também perguntou a cada participante a dificuldade percebida, em uma escala de 1 (muito fácil) a 5 (muito difícil), e se houve dúvidas sobre os requisitos apresentados. A dificuldade média declarada foi 3,4, e 9 dos 20 participantes (45,0\%) relataram dúvidas. Essas duas medidas não se concentram nos participantes cujos \textit{templates} reprovaram mais etapas da execução funcional, discutida na Seção~\ref{subsec:localstack}: dos 3 participantes com o pior resultado funcional (aprovação em apenas 1 dos 3 recursos), todos relataram dificuldade entre 3 e 4 e nenhum relatou dúvidas, valores próximos da média geral da amostra.

\subsection{Completude e Padronização dos Artefatos}

O \textit{workflow} \texttt{metrics.yml} inventariou os 240 \textit{templates} avaliados, verificando a presença dos arquivos esperados. Os quatro arquivos Terraform (\texttt{main.tf}, \texttt{variables.tf}, \texttt{outputs.tf} e \texttt{versions.tf}) estavam presentes nos 180 \textit{templates} automatizados e nos 60 \textit{templates} manuais (100\%).

O \texttt{README.md} esteve presente nos 180 \textit{templates} automatizados e em 42 dos 60 \textit{templates} manuais (70,0\%, 14 dos 20 participantes com documentação nos três recursos), embora a documentação constasse dos requisitos gerais fornecidos aos participantes. 

\subsection{Validação Sintática e Estrutural}

A Tabela~\ref{tab:terraform-validate} apresenta a taxa de aprovação em cada etapa do \textit{workflow} \texttt{terraform-validate.yml}. O cenário manual apresenta a maior aprovação geral (93,3\%, 56 de 60), seguido pelo cenário com contexto (88,9\%, 80 de 90) e pelo cenário sem contexto (62,2\%, 56 de 90).

\begin{table}[H]
\centering
\caption{Resultado do \texttt{terraform-validate} por cenário e etapa. Fonte: elaborado pelo autor.}
\label{tab:terraform-validate}
\small
\begin{tabular}{l|c|c|c|c}
\hline
\textbf{Cenário} & \textbf{fmt} & \textbf{init} & \textbf{validate*} & \textbf{Geral} \\
\hline
Manual (n=60) & 56/60 (93,3\%) & 60/60 (100\%) & 60/60 (100\%) & 56/60 (93,3\%) \\
IA sem contexto (n=90) & 69/90 (76,7\%) & 89/90 (98,9\%) & 76/89 (85,4\%) & 56/90 (62,2\%) \\
IA com contexto (n=90) & 82/90 (91,1\%) & 88/90 (97,8\%) & 88/88 (100\%) & 80/90 (88,9\%) \\
\hline
\end{tabular}
\vspace{2pt}
\footnotesize{*A coluna \texttt{validate} considera apenas os \textit{templates} que passaram pela etapa \texttt{init}, uma vez que a falha nessa etapa impede a execução da validação subsequente.}
\end{table}

As diferenças entre os cenários automatizados se concentram nas etapas de formatação (\texttt{fmt}) e de verificação de código (\texttt{validate}), enquanto a etapa de inicialização (\texttt{init}) aprova quase todos os \textit{templates} nos três cenários (detalhamento por recurso no Apêndice~\ref{apx:validacao}). No cenário sem contexto, a maior parte das reprovações vem de divergências simples de formatação (21 dos 90 \textit{templates}, 14 delas no IAM) e de um grupo menor de erros de escrita do código, como referência a uma variável inexistente, um trecho mal formado ao expor o algoritmo de criptografia do S3 e um valor inválido dentro de uma declaração repetida na política de confiança do IAM. No cenário com contexto, essas falhas praticamente desaparecem (8,9\% em \texttt{fmt} e nenhuma em \texttt{validate}), e as duas únicas reprovações de inicialização decorrem do mesmo provedor declarado duas vezes no mesmo arquivo. No cenário manual, as quatro reprovações são exclusivamente de formatação, sem nenhum erro de inicialização ou de verificação de código. Como uma divergência apenas de formatação já reprova o \textit{template} na coluna Geral, parte da diferença entre os cenários é de forma, e não de conteúdo. Desconsiderando essa etapa, a aprovação seria de 84,4\% (76/90) sem contexto, 97,8\% (88/90) com contexto e 100\% (60/60) no manual.

\subsection{Execução Funcional com LocalStack}
\label{subsec:localstack}

A Tabela~\ref{tab:localstack} apresenta a taxa de sucesso da execução funcional (\texttt{init}, \texttt{plan}, \texttt{apply}, verificação e \texttt{destroy}) no \textit{workflow} \texttt{localstack.yml}, por recurso e cenário.

\begin{table}[H]
\centering
\caption{Resultado do teste funcional com LocalStack por recurso e cenário. Fonte: elaborado pelo autor.}
\label{tab:localstack}
\small
\begin{tabular}{l|c|c|c}
\hline
\textbf{Recurso} & \textbf{Manual (n=20)} & \textbf{IA sem contexto (n=30)} & \textbf{IA com contexto (n=30)} \\
\hline
Amazon S3 & 12/20 (60,0\%) & 26/30 (86,7\%) & 29/30 (96,7\%) \\
AWS IAM & 19/20 (95,0\%) & 27/30 (90,0\%) & 30/30 (100\%) \\
Security Group & 16/20 (80,0\%) & 29/30 (96,7\%) & 27/30 (90,0\%) \\
\hline
Geral & 47/60 (78,3\%) & 82/90 (91,1\%) & 86/90 (95,6\%) \\
\hline
\end{tabular}
\end{table}

As 25 falhas foram classificadas pela mensagem de erro registrada nos \textit{logs} de cada etapa, com detalhamento no Apêndice~\ref{apx:localstack}. A maior parte das reprovações no cenário manual tem origem no próprio ambiente de teste, e não em um defeito do \textit{template}. A causa mais comum é uma operação que o \textit{provider} da AWS tenta executar sobre o recurso recém-criado, como consultar as \textit{tags} do \textit{bucket} ou gerar uma chave de criptografia, e que o LocalStack não oferece, responsável por 8 das 13 reprovações manuais de S3. As demais falhas manuais decorrem de uma variável obrigatória sem valor no IAM e, no Security Group, de descrições de regra escritas com acentuação (rejeitadas pela AWS) e de uma rotina de verificação que não reconhece o identificador do recurso criado por usar um nome ou formato diferente do esperado. Classificando essas causas entre o que é atribuível ao ambiente de execução e o que é atribuível ao conteúdo do \textit{template} (Apêndice~\ref{apx:localstack}), 11 das 13 falhas manuais (84,6\%) vêm do ambiente, e apenas 2 (15,4\%) refletem um defeito do próprio \textit{template}.

Cruzando o resultado funcional de cada participante manual com o formulário de caracterização, a relação com a experiência prévia em IaC é fraca nesta amostra. Dos participantes que já haviam utilizado alguma ferramenta de IaC antes do experimento, metade aprovou os três recursos, contra 40\% entre os que nunca haviam utilizado nenhuma. Os três participantes com o pior resultado (aprovação em apenas 1 dos 3 recursos) incluem dois sem nenhuma experiência prévia e um com conhecimento intermediário declarado e uso frequente de ferramentas de IA, e todos reprovam pelas mesmas causas ambientais listadas acima, não por um erro de conteúdo específico do \textit{template}.

Nos cenários de IA, as poucas falhas vêm mais de erros de escrita do código do que de limitações do ambiente. No cenário sem contexto, dois \textit{templates} erram a sintaxe de uma referência (no S3 e no IAM) e um terceiro declara um recurso experimental no Security Group. No cenário com contexto, as quatro falhas se dividem entre um recurso experimental, o mesmo provedor declarado duas vezes e dois \textit{templates} de Security Group cuja regra de validação da variável de nome rejeita o valor que o ambiente de execução fornece.

\subsection{Análise de Segurança}

Os \textit{templates} dos três cenários foram submetidos ao \textit{workflow} \texttt{security-scan.yml}, que executa Checkov e Trivy de forma independente sobre os arquivos Terraform. Como os achados dessas ferramentas compõem os próprios dados da pesquisa, o \textit{workflow} não interrompe a execução ao encontrá-los: uma execução concluída indica apenas que a análise foi realizada, não que os \textit{templates} estão livres de problemas.

\subsubsection{Configurações Inseguras (Trivy)}

A Tabela~\ref{tab:trivy} apresenta a proporção de \textit{templates} com pelo menos um achado do Trivy, por recurso, severidade e cenário.

\begin{table}[H]
\centering
\caption{Resultado da análise de configurações inseguras com Trivy, por recurso, severidade e cenário. Fonte: elaborado pelo autor.}
\label{tab:trivy}
\begin{tabular}{l|l|c|c|c}
\hline
\textbf{Recurso} & \textbf{Severidade} & \textbf{Manual} & \textbf{IA sem contexto} & \textbf{IA com contexto} \\
\hline
Amazon S3 & Baixa & 20/20 (100\%) & 30/30 (100\%) & 30/30 (100\%) \\
Amazon S3 & Média & 1/20 (5,0\%) & 0/30 (0,0\%) & 0/30 (0,0\%) \\
Amazon S3 & Alta & 17/20 (85,0\%) & 0/30 (0,0\%) & 21/30 (70,0\%) \\
\hline
Security Group & Baixa & 1/20 (5,0\%) & 0/30 (0,0\%) & 1/30 (3,3\%) \\
Security Group & Crítica & 9/20 (45,0\%) & 28/30 (93,3\%) & 8/30 (26,7\%) \\
\hline
AWS IAM & --- & 0/20 (0,0\%) & 0/30 (0,0\%) & 0/30 (0,0\%) \\
\hline
\end{tabular}
\vspace{2pt}
\footnotesize{AWS IAM não apresenta achados de nenhuma severidade em nenhum dos três cenários.}
\end{table}

O achado mais relevante é uma regra de severidade crítica que sinaliza quando o Security Group libera a porta de saída para qualquer endereço da internet, presente em 28 de 30 \textit{templates} (93,3\%) sem contexto, ainda que o \textit{prompt} desse cenário pedisse regras de saída explícitas, evitando abertura irrestrita desnecessária. No cenário com contexto, cuja regra organizacional proíbe esse endereço fora da porta 443/tcp, a incidência cai para 26,7\%, patamar próximo ao do manual (45,0\%, 9 de 20). Um achado de severidade baixa relativo à ausência de descrição da regra ocorre em um \textit{template} manual e em um com contexto.

No S3, os três cenários apresentam, em praticamente todos os \textit{templates}, um achado de severidade baixa relativo à ausência de registro de acesso, controle que não consta dos requisitos de nenhum dos grupos. Um achado de severidade alta relativo à ausência de chave de criptografia gerenciada pelo cliente ocorre em 17 de 20 \textit{templates} manuais e em 21 de 30 com contexto, mas em nenhum sem contexto. Os \textit{templates} sem contexto deixam essa chave como uma opção configurável por variável, usando a criptografia padrão da AWS quando nenhuma é informada, o que a análise estática do Trivy não sinaliza, enquanto o contexto organizacional estabelece essa criptografia padrão como exigência e a maioria dos participantes manuais adota a mesma escolha. Um único achado de severidade média, relativo à ausência de versionamento, ocorre em um \textit{template} manual.

\subsubsection{Conformidade com Boas Práticas (Checkov)}

O Checkov avaliou 4.559 verificações sobre os 240 \textit{templates}, sendo 1.183 no cenário manual, 1.585 no cenário sem contexto e 1.791 no cenário com contexto. A aprovação geral foi de 89,3\% no manual, 89,6\% sem contexto e 89,9\% com contexto, valores muito próximos entre si.

Essa proximidade decorre da distribuição dos achados por recurso (Tabelas~\ref{tab:apx-checkov-recurso} e~\ref{tab:apx-checkov-regra} do Apêndice~\ref{apx:seguranca}). Como o número de verificações aplicáveis varia por recurso, o agregado acima pesa mais os \textit{templates} com mais verificações, e o S3, que concentra a maioria das reprovações, puxa a média geral para baixo. Calculando a taxa de aprovação de cada \textit{template} individualmente e depois tirando a média entre eles, o resultado muda pouco (87,4\% no manual, 86,8\% sem contexto e 89,9\% com contexto, mediana de 93,3\%, 87,5\% e 87,5\%, respectivamente, com detalhamento no Apêndice~\ref{apx:seguranca}). Nos três cenários, todos os \textit{templates} de IAM aprovam a totalidade das verificações aplicáveis, e praticamente todos os de Security Group reprovam a mesma verificação, que exige que o grupo esteja associado a outro recurso da infraestrutura e que nenhum deles, avaliado isoladamente, atende. A diferença está concentrada no S3, com 75,8\% de aprovação no manual, 72,6\% sem contexto e 81,7\% com contexto, em controles que não constam dos requisitos de nenhum dos três grupos, como ciclo de vida dos objetos, notificação de eventos, replicação entre regiões e registro de acesso.

Em síntese, o contexto organizacional faz mais diferença no risco de exposição de rede do Security Group do que na conformidade ampla medida pelo Checkov, na qual as diferenças entre os três cenários (89,3\% a 89,9\%) são numericamente pequenas na amostra observada.

\section{Ameaças à Validade}
\label{sec:ameacas}

A principal limitação é o tamanho da amostra manual. Foram 60 \textit{templates} de 20 participantes, contra 90 por cenário de IA. Esse tamanho foi limitado pela dificuldade de recrutar voluntários que concluíssem a atividade dentro do período de coleta, não por qualquer exigência de experiência prévia, que não houve. Terraform é uma tecnologia pouco presente nas disciplinas de graduação e no dia a dia da maioria dos estudantes de tecnologia, o que reduziu o número de pessoas dispostas e disponíveis para terminar a atividade manual no prazo. Por causa disso, os resultados do cenário manual são mais sensíveis a casos isolados do que os dos cenários automatizados, que têm 30 execuções cada.

Há também limitações no ambiente de teste. O LocalStack usa uma lista de variáveis de entrada e um roteiro de verificação por \textit{output}, os dois montados a partir dos nomes mais comuns encontrados nos \textit{templates}. Um \textit{template} que usa um nome de variável ou de \textit{output} fora dessa lista, ou que expõe o \textit{output} em um formato diferente, como um mapa com vários identificadores em vez de um único valor, pode não ser testado corretamente nessa etapa, mesmo que o recurso em si esteja certo. O emulador também não implementa os serviços S3 Control e KMS da AWS. Então, \textit{templates} cujo código usa essas APIs, direta ou indiretamente, reprovam a execução funcional por uma limitação do ambiente, não por um erro do \textit{template}. Por fim, a verificação \texttt{CKV2\_AWS\_5} do Checkov sempre reprova nos três cenários, porque ela exige que o Security Group esteja associado a outro recurso, e isso é esperado no desenho deste experimento, que avalia cada Security Group isoladamente.

A padronização foi medida apenas pela presença dos arquivos esperados, sem verificar a aderência às convenções de nomenclatura e \textit{tags} do contexto organizacional. Três participantes manuais enviaram o próprio arquivo de travamento de dependências, e seus \textit{templates} foram executados no LocalStack com o \textit{provider} AWS 6.68.0, e não com o 5.100.0 usado nos demais.

Os resultados refletem um único modelo (Claude, variante \textit{Sonnet}, identificado nos metadados só por esse nome, sem a versão exata) e um único conjunto de \textit{prompts} e documento de contexto organizacional. Não é objetivo deste trabalho comparar modelos de IA entre si, mas sim comparar métodos de criação de \textit{templates}. Por isso, testar outros modelos, versões ou formulações de contexto fica como trabalho futuro.

\section{Conclusão}
\label{sec:conclusao}

O primeiro objetivo específico era analisar se a IA reduz o tempo e o esforço de criar os \textit{templates}. Os dados confirmam isso de forma direta: a geração automatizada leva 2,7 minutos sem contexto e 3,3 minutos com contexto, contra 77,5 minutos de desenvolvimento manual líquido. A diferença entre IA e manual é muito maior do que o custo de acrescentar contexto organizacional, o que sugere que fornecer contexto é uma forma barata de buscar mais qualidade sem perder o ganho de tempo. Como discutido na Seção~\ref{sec:resultados}, a dificuldade percebida e as dúvidas relatadas pelos participantes também não se concentram em quem teve o \textit{template} reprovado na execução funcional, o que indica que o quanto a pessoa achou a tarefa difícil não é um bom indício da qualidade final do artefato.

O segundo objetivo era avaliar a conformidade dos \textit{templates} com boas práticas e padrões de IaC, e aqui o resultado depende da etapa observada. Na validação sintática, o manual tem a maior aprovação (93,3\%), seguido pelo com contexto (88,9\%) e pelo sem contexto (62,2\%). Já na execução funcional com LocalStack essa ordem se inverte: o com contexto tem a maior aprovação (95,6\%), seguido pelo sem contexto (91,1\%) e pelo manual (78,3\%). Isso mostra que o contexto organizacional ajuda sobretudo a IA a seguir convenções de formatação, o tipo de regra que um \textit{prompt} mais específico resolve com facilidade. O manual, por outro lado, tende a escrever código sintaticamente correto, mas tropeça mais na execução real: oito das treze reprovações do cenário manual vêm da mesma causa ambiental, quando o \textit{provider} da AWS tenta usar uma operação que o ambiente de teste não suporta, e isso não diz nada sobre a qualidade do \textit{template} em si. Como também discutido na Seção~\ref{sec:resultados}, a relação entre experiência prévia em IaC e aprovação funcional é fraca nesta amostra, e os piores resultados reprovam pelas mesmas causas ambientais que afetam qualquer nível de experiência. Isso sugere que o conhecimento prévio em Terraform conta menos do que se esperava para produzir um \textit{template} funcional, pelo menos nesta amostra. A conformidade ampla medida pelo Checkov não acompanha essa diferença: os três cenários ficam muito próximos (89,3\% a 89,9\%), porque a maior parte das reprovações recai sobre controles que não fazem parte dos requisitos de nenhum grupo, como ciclo de vida e replicação do S3.

O terceiro objetivo era investigar vulnerabilidades de segurança produzidas pelas diferentes formas de geração. A diferença mais clara aparece na saída irrestrita do Security Group, presente em 93,3\% dos \textit{templates} sem contexto, contra 26,7\% com contexto e 45,0\% no manual. Isso mostra que pedir, de forma genérica, para evitar abertura irrestrita não é suficiente para a IA sem contexto, enquanto uma regra organizacional concreta, como proibir \texttt{0.0.0.0/0} fora da porta 443, aproxima o resultado da IA do manual. No S3, a ausência de chave gerenciada pelo cliente aparece no manual (85,0\%) e no com contexto (70,0\%), mas não no sem contexto. Isso não é uma falha de segurança introduzida pela IA, e sim reflexo de uma escolha de criptografia (SSE-S3) que o manual e o cenário com contexto adotaram e o sem contexto não.

No geral, a IA generativa produz \textit{templates} funcionais em uma fração do tempo humano, e apresenta resultados comparáveis ou superiores ao desenvolvimento manual em parte das métricas avaliadas, sobretudo na execução funcional e na redução de configurações inseguras de rede, quando recebe um contexto organizacional explícito. Esse mesmo contexto é o que mais reduz o risco de exposição de rede. Nesta amostra, os resultados sugerem que a qualidade dos \textit{templates}, manuais ou gerados por IA, esteve mais associada a regras explícitas e verificáveis, como nomes de variáveis, restrições de rede e padrões de formatação, do que à experiência prévia de quem os produziu ou à ausência de orientação adicional dada à IA. Por isso, recomenda-se que a adoção de IA na geração de IaC venha acompanhada de um contexto organizacional explícito e da mesma validação automatizada usada no desenvolvimento manual, desenhada para reconhecer variações legítimas de nome e de formato de \textit{output} sem confundi-las com defeito do recurso avaliado.

Como trabalhos futuros, propõe-se um ciclo totalmente agêntico, em que o retorno dos \textit{pipelines} volta para o modelo corrigir o \textit{template} até passar em todas as verificações. Propõe-se também ampliar os cenários para outros recursos da AWS, como EC2, RDS, Lambda e VPC, incluindo arquiteturas com dependência entre recursos, e expressar o contexto organizacional como políticas verificáveis por código. Por fim, um cenário híbrido, com revisão humana sobre \textit{templates} gerados com contexto, e a repetição do experimento com outros modelos e ferramentas de IaC, como OpenTofu e AWS CloudFormation, ajudariam a confirmar o alcance desses resultados.


\vspace{-2mm}
\section*{Pacote de Replicação}
\vspace{-1mm}
\footnotesize
O pacote de replicação deste trabalho encontra-se disponível em\\
\url{https://github.com/ICEI-PUC-Minas-PPLES-TI/}\\
\url{plf-es-2026-1-tcci-3585100-pes-david-leong}
\normalsize

\FloatBarrier

\bibliographystyle{sbc}
\bibliography{sbc-template}

\section*{Apêndice}
\appendix
\section{Validação Sintática e Completude por Recurso}
\label{apx:validacao}

A Tabela~\ref{tab:apx-validate} apresenta o resultado de cada etapa do \texttt{terraform-validate.yml} por cenário e recurso. Como na Tabela~\ref{tab:terraform-validate}, a coluna \texttt{validate} considera apenas os \textit{templates} aprovados no \texttt{init}.

\begin{table}[H]
\centering
\caption{Resultado do \texttt{terraform-validate} por cenário, recurso e etapa. Fonte: elaborado pelo autor.}
\label{tab:apx-validate}
\small
\begin{tabular}{l|l|c|c|c|c}
\hline
\textbf{Cenário} & \textbf{Recurso} & \textbf{fmt} & \textbf{init} & \textbf{validate} & \textbf{Geral} \\
\hline
Manual & Amazon S3 & 20/20 & 20/20 & 20/20 & 20/20 (100\%) \\
 & AWS IAM & 19/20 & 20/20 & 20/20 & 19/20 (95,0\%) \\
 & Security Group & 17/20 & 20/20 & 20/20 & 17/20 (85,0\%) \\
\hline
IA sem contexto & Amazon S3 & 28/30 & 30/30 & 23/30 & 21/30 (70,0\%) \\
 & AWS IAM & 16/30 & 30/30 & 27/30 & 14/30 (46,7\%) \\
 & Security Group & 25/30 & 29/30 & 26/29 & 21/30 (70,0\%) \\
\hline
IA com contexto & Amazon S3 & 29/30 & 29/30 & 29/29 & 28/30 (93,3\%) \\
 & AWS IAM & 24/30 & 30/30 & 30/30 & 24/30 (80,0\%) \\
 & Security Group & 29/30 & 29/30 & 29/29 & 28/30 (93,3\%) \\
\hline
\end{tabular}
\end{table}

No cenário manual, as quatro reprovações em \texttt{fmt} ocorrem três no Security Group e uma no IAM, por divergências de formatação. O IAM concentra a maior parte das reprovações em \texttt{fmt} nos dois cenários de IA (14 de 21 sem contexto e 6 de 8 com contexto). As reprovações em \texttt{validate} do cenário sem contexto ocorrem no S3 (7), no IAM (3) e no Security Group (3). As falhas de \texttt{init} decorrem de recurso experimental declarado no bloco \texttt{terraform} (1 sem contexto) e de declaração duplicada de \texttt{required\_providers} (2 com contexto).

Quanto à completude, os quatro arquivos Terraform (\texttt{main.tf}, \texttt{variables.tf}, \texttt{outputs.tf} e \texttt{versions.tf}) estavam presentes nos 180 \textit{templates} automatizados e nos 60 \textit{templates} manuais (100\%). O \texttt{README.md} esteve presente em 30 de 30 \textit{templates} de cada recurso nos dois cenários de IA e em 42 dos 60 \textit{templates} manuais, de 14 dos 20 participantes.

\section{Causas das Falhas na Execução Funcional}
\label{apx:localstack}

A Tabela~\ref{tab:localstack-causas} classifica as 25 falhas da execução funcional com LocalStack pela mensagem de erro registrada nos \textit{logs} de cada etapa.

\begin{table}[H]
\centering
\caption{Causas das falhas na execução funcional com LocalStack, por cenário. Fonte: elaborado pelo autor.}
\label{tab:localstack-causas}
\small
\begin{tabularx}{\textwidth}{>{\raggedright\arraybackslash}X|c|c|c}
\hline
\textbf{Causa (mensagem do Terraform ou do \textit{workflow})} & \textbf{Manual} & \textbf{IA s/ contexto} & \textbf{IA c/ contexto} \\
\hline
Operação não suportada pelo ambiente de execução (S3 Control ou KMS) & 7 & 0 & 0 \\
Conflito de nome com outro \textit{template} executado na mesma rodada & 1 & 0 & 0 \\
Variável obrigatória sem valor no ambiente de execução & 1 & 0 & 0 \\
Descrição de regra do Security Group com acentos, rejeitada pelo \textit{provider} & 2 & 0 & 0 \\
\textit{Output} em formato ou nome não reconhecido pela rotina de verificação & 2 & 0 & 0 \\
Indexação de conjunto em \textit{output} do S3 & 0 & 4 & 0 \\
\texttt{for\_each} dinâmico inválido na \textit{trust policy} do IAM & 0 & 3 & 0 \\
Recurso experimental declarado no bloco \texttt{terraform} & 0 & 1 & 1 \\
Declaração duplicada de \texttt{required\_providers} & 0 & 0 & 1 \\
Valor do ambiente de execução rejeitado pela \texttt{validation} do \textit{template} & 0 & 0 & 2 \\
\hline
\textbf{Total} & \textbf{13 de 60} & \textbf{8 de 90} & \textbf{4 de 90} \\
\hline
\end{tabularx}
\end{table}

A Tabela~\ref{tab:localstack-atribuicao} classifica as 13 falhas manuais entre as atribuíveis ao ambiente de execução (nomes de variável e \textit{output} fora do repertório reconhecido, e operações não suportadas pelo emulador) e as atribuíveis ao conteúdo do \textit{template} (um erro que também ocorreria em uma conta AWS real).

\begin{table}[H]
\centering
\caption{Falhas manuais no LocalStack, por origem atribuída. Fonte: elaborado pelo autor.}
\label{tab:localstack-atribuicao}
\small
\begin{tabularx}{\textwidth}{>{\raggedright\arraybackslash}X|c|c}
\hline
\textbf{Origem} & \textbf{Falhas} & \textbf{\%} \\
\hline
Ambiente de execução (S3 Control/KMS, conflito de nome, variável e \textit{output} fora do repertório) & 11 & 84,6\% \\
\hline
Conteúdo do \textit{template} (descrição de regra com acentuação, rejeitada pela própria AWS) & 2 & 15,4\% \\
\hline
Total & 13 & 100\% \\
\hline
\end{tabularx}
\end{table}

No cenário manual, as oito reprovações de S3 concentram-se em uma única causa. O \textit{provider} AWS invoca, ao criar o \textit{bucket} ou a chave de criptografia, uma operação de listagem de \textit{tags} via S3 Control ou de criação de chave via KMS, serviços que o ambiente de execução local não implementa. Em um dos casos, o \textit{bucket} de um \textit{template} anterior, executado na mesma rodada, permanece no ambiente e provoca um conflito de nome em outro \textit{template} que declara o mesmo valor fixo. As demais cinco reprovações, no IAM e no Security Group, têm três causas distintas entre si: uma variável obrigatória sem valor, duas descrições de regra em português com acentuação e dois casos em que o recurso é criado corretamente, mas a rotina de verificação não localiza seu identificador, por o \textit{output} estar em um nome fora do padrão reconhecido ou em um mapa indexado por múltiplos Security Groups do mesmo \textit{template}.

No cenário sem contexto, as falhas exclusivas desse cenário (indexação de conjunto no S3 e \texttt{for\_each} inválido no IAM) são erros de construção da linguagem, também detectados pelo \texttt{validate}. No cenário com contexto, as reprovações do Security Group decorrem de dois \textit{templates} que transformaram o padrão de nomenclatura organizacional em regra de \texttt{validation} da variável de nome e, por isso, rejeitaram o valor fornecido pelo ambiente de execução; a reprovação do S3 decorre de uma declaração duplicada de \texttt{required\_providers}.

\section{Detalhes da Análise de Segurança}
\label{apx:seguranca}

A Tabela~\ref{tab:apx-checkov-recurso} apresenta o número de verificações do Checkov e a taxa de aprovação por recurso e cenário, a Tabela~\ref{tab:apx-checkov-regra} apresenta, para cada verificação reprovada, o número de \textit{templates} com pelo menos uma reprovação, e a Tabela~\ref{tab:apx-checkov-template} apresenta a taxa de aprovação calculada por \textit{template} individual, em vez do agregado de verificações, como alternativa que não pesa mais os recursos com mais verificações aplicáveis.

\begin{table}[H]
\centering
\caption{Taxa de aprovação do Checkov calculada por \textit{template} individual, por cenário. Fonte: elaborado pelo autor.}
\label{tab:apx-checkov-template}
\small
\begin{tabular}{l|c|c|c|c}
\hline
\textbf{Cenário} & \textbf{Média} & \textbf{Mediana} & \textbf{Mínimo} & \textbf{Máximo} \\
\hline
Manual & 87,4\% & 93,3\% & 41,7\% & 100\% \\
IA sem contexto & 86,8\% & 87,5\% & 66,7\% & 100\% \\
IA com contexto & 89,9\% & 87,5\% & 70,6\% & 100\% \\
\hline
\end{tabular}
\end{table}

\begin{table}[H]
\centering
\caption{Verificações do Checkov e taxa de aprovação por recurso e cenário. Fonte: elaborado pelo autor.}
\label{tab:apx-checkov-recurso}
\small
\begin{tabular}{l|c|c|c}
\hline
\textbf{Recurso} & \textbf{Manual} & \textbf{IA sem contexto} & \textbf{IA com contexto} \\
\hline
Amazon S3 & 317/418 (75,8\%) & 357/492 (72,6\%) & 668/818 (81,7\%) \\
AWS IAM & 471/471 (100\%) & 836/836 (100\%) & 705/705 (100\%) \\
Security Group & 269/294 (91,5\%) & 227/257 (88,3\%) & 238/268 (88,8\%) \\
\hline
Total & 1.057/1.183 (89,3\%) & 1.420/1.585 (89,6\%) & 1.611/1.791 (89,9\%) \\
\hline
\end{tabular}
\end{table}

\begin{table}[H]
\centering
\caption{\textit{Templates} com pelo menos uma reprovação em cada verificação do Checkov, por cenário. Fonte: elaborado pelo autor.}
\label{tab:apx-checkov-regra}
\small
\begin{tabularx}{\textwidth}{l|>{\raggedright\arraybackslash}X|c|c|c}
\hline
\textbf{Verificação} & \textbf{Descrição} & \textbf{Manual} & \textbf{IA s/ ctx.} & \textbf{IA c/ ctx.} \\
\hline
\texttt{CKV2\_AWS\_61} & S3 sem configuração de ciclo de vida & 20/20 & 30/30 & 30/30 \\
\texttt{CKV2\_AWS\_62} & S3 sem notificação de eventos & 20/20 & 30/30 & 30/30 \\
\texttt{CKV\_AWS\_144} & S3 sem replicação entre regiões & 20/20 & 30/30 & 30/30 \\
\texttt{CKV\_AWS\_18} & S3 sem \textit{log} de acesso & 19/20 & 30/30 & 30/30 \\
\texttt{CKV\_AWS\_145} & S3 sem criptografia KMS por padrão & 15/20 & 15/30 & 30/30 \\
\texttt{CKV2\_AWS\_64} & Chave KMS sem política definida & 2/20 & 0/30 & 0/30 \\
\texttt{CKV2\_AWS\_6} & S3 sem bloqueio de acesso público & 1/20 & 0/30 & 0/30 \\
\texttt{CKV\_AWS\_21} & S3 sem versionamento & 1/20 & 0/30 & 0/30 \\
\hline
\texttt{CKV2\_AWS\_5} & Security Group não associado a outro recurso & 20/20 & 30/30 & 30/30 \\
\texttt{CKV\_AWS\_23} & Security Group com regra sem descrição & 1/20 & 0/30 & 0/30 \\
\texttt{CKV\_AWS\_382} & Security Group com saída irrestrita & 2/20 & 0/30 & 0/30 \\
\hline
\end{tabularx}
\end{table}

A Tabela~\ref{tab:apx-trivy-regra} apresenta, para cada regra do Trivy com achados, o número de \textit{templates} afetados e o total de ocorrências.

\begin{table}[H]
\centering
\caption{Achados do Trivy por regra e cenário: \textit{templates} afetados (ocorrências). Fonte: elaborado pelo autor.}
\label{tab:apx-trivy-regra}
\small
\begin{tabularx}{\textwidth}{l|>{\raggedright\arraybackslash}X|c|c|c|c}
\hline
\textbf{Regra} & \textbf{Descrição} & \textbf{Sev.} & \textbf{Manual} & \textbf{IA s/ ctx.} & \textbf{IA c/ ctx.} \\
\hline
\texttt{AWS-0089} & S3 sem \textit{log} de acesso & Baixa & 20 (20) & 30 (30) & 30 (30) \\
\texttt{AWS-0090} & S3 sem versionamento & Média & 1 (1) & 0 (0) & 0 (0) \\
\texttt{AWS-0132} & S3 sem chave gerenciada pelo cliente & Alta & 17 (17) & 0 (0) & 21 (21) \\
\texttt{AWS-0104} & Saída irrestrita no Security Group & Crítica & 9 (10) & 28 (28) & 8 (8) \\
\texttt{AWS-0099} & Security Group sem descrição & Baixa & 1 (1) & 0 (0) & 1 (1) \\
\hline
\end{tabularx}
\end{table}


\end{document}
